% Generated draft: requires mathematical and source review.
\documentclass[11pt]{article}
\usepackage[margin=25mm]{geometry}
\usepackage{fontspec}
\setmainfont{Latin Modern Roman}
\usepackage{amsmath,amssymb,mathtools,unicode-math}
\setmathfont{Latin Modern Math}
\usepackage{enumitem,needspace,xurl,hyperref}
\setlength{\parindent}{0pt}
\setlength{\parskip}{5pt}
\setlength{\emergencystretch}{3em}
\newcommand{\N}{\mathbb{N}}
\newcommand{\Z}{\mathbb{Z}}
\newcommand{\Q}{\mathbb{Q}}
\newcommand{\R}{\mathbb{R}}
\newcommand{\C}{\mathbb{C}}
\newcommand{\F}{\mathbb{F}}
\newcommand{\A}{\mathbb{A}}
\newcommand{\PP}{\mathbb{P}}
\newcommand{\E}{\mathbb{E}}
\newcommand{\eps}{\varepsilon}
\newcommand{\dd}{\,\mathrm{d}}
\begin{document}
\section*{UnsolvedMath ID: 30000472}
\textbf{Source problem code:} catalogue entry
\subsection{Definitions and mathematical statement}
For a regular uncountable cardinal $\mu$, a $\mu$-Aronszajn tree has height $\mu$, levels of size less than $\mu$, and no branch of order type $\mu$. At $\mu=\nu^+$, it is special if its nodes are the union of $\nu$ antichains, equivalently if a map to $\nu$ is injective on each chain. Absence of all Aronszajn trees is the tree property and is stronger than absence of special Aronszajn trees.

The restored Question 6 asks for a particular forcing method, not merely a consistency conclusion. Start in a model with cardinals $\kappa<\lambda$, with $\kappa$ $\lambda$-supercompact and $\lambda$ weakly compact. In the standard embedding interpretation, the former gives an elementary embedding with critical point $\kappa$, image of $\kappa$ above $\lambda$, and target closed under $\lambda$-sequences; weak compactness has its usual large-cardinal meaning. Can an iteration of the report's nonstationary-ideal forcing method reprove the cited Abraham conclusion, producing a generic extension $V[G]$ in which
\[
 \text{there are no }\aleph_2^{V[G]}\text{-Aronszajn trees and no }
                 \aleph_3^{V[G]}\text{-Aronszajn trees}?
\]
The cardinal symbols in the target refer to the extension. The source does not specify the collapse bookkeeping identifying ground-model cardinals with them.

As a weaker fallback, it asks whether that method can at least construct a model with no special Aronszajn trees on $\aleph_2$ or $\aleph_3$. The source's ``or'' is retained: the precise requirement of one target versus both simultaneously needs confirmation. This fallback must not be strengthened to the full tree property without proof.

The particular nonstationary-ideal forcing posets, iteration support, preservation arguments and any extra hypotheses are not supplied in the extract. A construction by another technique does not answer the method-specific request. The literature assessment checked on 21 August 2026 reports consistency models with no special trees at many successive regular cardinals, including $\aleph_2$, but does not claim the requested method reproves Abraham's theorem from exactly the displayed assumptions. These are dated supplied partial-progress claims requiring independent verification.
\subsection{Short English statement}
Can the report's nonstationary-ideal forcing iteration recover Abraham's simultaneous tree-property result at $\aleph_2$ and $\aleph_3$ under the stated large cardinals? The weaker fallback excludes only special trees, and existing constructions by other methods do not settle this request.
\subsection{Sources}
\begin{itemize}
\item \emph{Feinstrukturtheorie und Innere Modelle} (2007), Oberwolfach Reports, DOI 10.4171/owr/2006/20, pp. 1215--1244, Question 6 according to supplied verification, \url{https://doi.org/10.4171/owr/2006/20}; proposed year supplied as 2006. The restored theorem is attributed to Abraham; no separate full Abraham citation is supplied.
\item Literature assessment checked 21 August 2026: original report, \url{https://doi.org/10.4171/owr/2006/20}; Unger, \emph{Successive failures of approachability}, \url{https://arxiv.org/abs/1702.05062}.
\item Local UnsolvedMath registry, permanent ID $30000472$, exact source code \texttt{catalogue entry}, created 25 August 2026 and updated 16 September 2026. Supplied verification expands the referenced theorem and retains the method request.
\end{itemize}
\end{document}
